% !PRINTED 599
\weberpage{599}
%<<CONT
\weq{6}{\alpha_0 = \alpha + M\vartheta \equiv 0 \pmod{\mathfrak{m}}}
where $\vartheta$ is an integer in $\Omega$.

From this it follows that every prime ideal $\mathfrak{e}$ dividing the
excludent must divide $M$, since otherwise $\alpha_0$ would be divisible
by $\mathfrak{e}$ for $\mathfrak{m} = \mathfrak{e}$.

If $\alpha$ and $M$ have a common divisor, then all numbers of the form
$\alpha + M\xi$, in which $\xi$ is an arbitrary number, have the same
common divisor with $M$. But in $A$ there are also numbers that are
relatively prime to $M$, e.\,g. the number 1. The following theorem
then holds:

3. \gesperrt{If $\mathfrak{m}$ is relatively prime to $M$, then in $A$
one can determine an integer $\alpha_0 = \mathfrak{m}\mathfrak{n}$
divisible by $\mathfrak{m}$ in such a way that $\mathfrak{n}$ becomes
relatively prime to $M$ and to an arbitrarily given ideal $\mathfrak{c}$.}

For if by (6) one has determined a number divisible by $\mathfrak{m}$,
\[
\alpha_0 = \mathfrak{m}\mathfrak{n}
\]
which is coprime to $M$, then one obtains other such numbers:
\[
\alpha_0' = \alpha_0 + M\xi = \mathfrak{m}\mathfrak{n}',
\]
if
\[
\xi = \mathfrak{m}\mathfrak{r}
\]
is an arbitrary integer divisible by $\mathfrak{m}$. Now if $\mathfrak{q}$ is a prime ideal that divides $\mathfrak{c}$
but not $\alpha_0$, let $\mathfrak{r}$ be assumed divisible by
$\mathfrak{q}$; then $\mathfrak{n}'$ is not divisible by $\mathfrak{q}$. But if $\mathfrak{q}$ divides $\mathfrak{c}$
and $\alpha_0$, let $\mathfrak{r}$ be assumed not divisible by $\mathfrak{q}$; then
$\mathfrak{n}'$ is likewise not divisible by $\mathfrak{q}$.

From this it follows:

4. \gesperrt{If $\mathfrak{m}$ is an arbitrary ideal coprime to $M$,
then in every ideal class $\overline{A_\nu}$ there are ideals that are
relatively prime to $\mathfrak{m}$.}

To prove this, take an ideal $\mathfrak{b}$ in the class
$(\overline{A_\nu})^{-1}$ and determine by 3. a number
\[
\alpha = \mathfrak{a}\mathfrak{b},
\]
so that $\mathfrak{a}$ is relatively prime to $\mathfrak{m}$; then $\mathfrak{a}$ belongs to the class $\overline{A_\nu}$
and satisfies the requirement of Theorem 4.

\noindent\textbf{\S\ 162.\qquad \gesperrt{The ideals of the principal class divisible by an ideal.}}

\gesperrt{By $\Omega$ there shall now be understood the quadratic field with
negative fundamental number $\Delta$.}
